% Printed 62--71, including the intervening unnumbered leaves.
\clearpage
\noindent\makebox[2em][l]{62}\parbox{\dimexpr\linewidth-4em\relax}{\centering\small ANZEIGE.\quad BESTIMMUNG DES BREITENUNTERSCHIEDS ETC.}\hfill\mbox{}\par
\noindent Herr \textsc{Hansen}, nach Absprache mit dem Verfasser gleichzeitig mit diesem an
denselben Zenitsternen mit dem Meridiankreis von \textsc{Ertel} angestellt hat. Das auf
206 Beobachtungen gegründete Ergebnis gewinnt zusätzlich an Interesse, weil die Seeberger
Sternwarte durch die unter der Leitung von Generalleutnant von \textsc{Müffling} gemessenen
Dreiecke mit dem hannoverschen Dreieckssystem verbunden ist.

Zum Schluss wird noch das Ergebnis einer neuen Berechnung aller bisherigen zuverlässigen
Gradmessungen mitgeteilt, die Herr Dr.\ \textsc{Schmidt} auf Veranlassung des Verfassers
ausgeführt hat. Sie folgt ähnlichen Prinzipien wie die zuvor erwähnte Arbeit von \textsc{Walbeck},
rechnet aber genauer, berücksichtigt \textls[150]{sämtliche} in den einzelnen Gradmessungen
vorkommenden astronomischen Punkte (über die Endpunkte hinaus) und bezieht die Ergebnisse der
hannoverschen Gradmessung aus dem vorliegenden Werk ein. Diese Bestimmung der Dimensionen des
Erdellipsoids kann als die zuverlässigste Folgerung gelten, die sich bisher aus sämtlichen
Breitengradmessungen ziehen lässt. Es ist außerdem von Interesse zu sehen, wie sich unsere
Vorstellungen davon, welcher Grad an Genauigkeit auf diesem Weg bisher erreichbar ist, nach
bestimmten Prinzipien festlegen lassen.
\begin{center}\rule{.25\textwidth}{.4pt}\end{center}

\clearpage
% Unnumbered Bemerkungen opening at the position of printed 63.
\vspace*{14mm}
\begin{center}\textbf{BEMERKUNGEN.}\end{center}
{\small
Im Originaldruck des Breitenunterschieds, der auch anhand des noch erhaltenen Manuskripts von
\textsc{Gauss} überprüft werden konnte, fanden sich in den Zahlenangaben einige kleine
Ungenauigkeiten, die nur zum Teil berichtigt wurden.

Im vorstehenden Abdruck wurde auf S.\ 9 die Angabe des Originals geändert für
\begin{center}
\par\noindent\begin{tabular}{ccccccl}
die&Deklination&des&Sterns&26:&$50^\circ28'13''{,}12$&in $50^\circ38'13''{,}12$\\
»&»&»&»&40:&$56\quad42\quad21{,}84$&» $56\quad52\quad21{,}74$,
\end{tabular}
\end{center}
da sonst die auf S.\ 22 und 27 mitgeteilten Beobachtungswerte für diese Sterne mit den nach
der Ausgleichung verbleibenden Fehlern auf S.\ 35 nicht übereinstimmen.

Im II.\ Kapitel, den Beobachtungen, wurden vereinzelt kleine Änderungen von 1 bis 2 Hundertstel
Sekunden an den Endwerten vorgenommen, um sie mit den Beobachtungswerten und den Reduktionen
in Übereinstimmung zu bringen. Voraussetzung war, dass die in die spätere Rechnung eingehenden
mittleren Endwerte dadurch unverändert blieben. Die folgenden Stellen mussten unberichtigt bleiben.
\par\noindent\begin{tabular}{@{}p{.065\linewidth}p{.91\linewidth}@{}}
S.\ 21&Bei Stern 25 ($\theta$ Draconis), Altona. Teilkreis West. 16.\ Juni, ergibt sich aus dem
Beobachtungswert und den beiden Reduktionen ein um $1''$ größerer Endwert als angegeben.\\
S.\ 22&Bei Stern 26, Altona. Teilkreis Ost. 27.\ Juni, ergibt sich aus dem Beobachtungswert und den
Reduktionen der Endwert $\ldots\ 34''{,}98$ statt $\ldots\ 34''{,}89$. Der letztere Wert wurde
jedoch bei der Mittelbildung verwendet.\\
S.\ 26&Bei Stern 36, Göttingen. Teilkreis West, ist das Mittel der Endwerte $\ldots\ 21''{,}14$.
Für die Rechnung wurde der angegebene Wert $\ldots\ 21''{,}13$ verwendet.\\
S.\ 27&Bei Stern 40, Altona. Teilkreis West. 22.\ Juni, ist der aus dem Beobachtungswert und den
Reduktionen abgeleitete Endwert um $0''{,}10$ größer als dort angegeben.
\end{tabular}

Im III.\ Kapitel, Ergebnisse, wurden die folgenden Änderungen vorgenommen.
\par\noindent\begin{tabular}{@{}p{.065\linewidth}p{.91\linewidth}@{}}
S.\ 30&Der Breitenunterschied aus den Beobachtungen des Sterns 39, für den im Original
$2^\circ0'58''{,}51$ steht, wurde in $\ldots\ 58''{,}57$ geändert, wie sich aus den Werten von
S.\ 27 ergibt. Der angegebene Mittelwert für den Breitenunterschied auf S.\ 30 bleibt davon
unberührt. Das Gewicht des Ergebnisses ist, wie angegeben, $213{,}42$. Im Original steht
dafür $213{,}41$.\\
S.\ 34&In Art.\ 6, Z.\ 8 v.\ o., wurde die Angabe des Originals für die Korrektur des
Kollimationsfehlers in Göttingen von $-0''{,}012$ in $+0''{,}012$ geändert, entsprechend den
Werten $3''{,}75$ auf S.\ 32, Z.\ 2 v.\ o., und $3''{,}76$ auf S.\ 34, Z.\ 12 v.\ o.
Die Nachrechnung hat die Richtigkeit des $+$-Zeichens bestätigt.
\end{tabular}\par
}

\clearpage
\noindent\makebox[2em][l]{64}\parbox{\dimexpr\linewidth-4em\relax}{\centering\small BEMERKUNGEN.\quad BESTIMMUNG DES BREITENUNTERSCHIEDS ETC.}\hfill\mbox{}\par
{\small
\par\noindent\begin{tabular}{@{}p{.065\linewidth}p{.91\linewidth}@{}}
S.\ 45&In der Fehlerspalte der Tabelle
\end{tabular}
\begin{center}
\par\noindent\begin{tabular}{llllll}
der erste&Wert&$-0''{,}31$&des Originals&in&$-0''{,}30$,\\
der letzte&»&$-0{,}23$&»\quad »&»&$-0{,}24$;
\end{tabular}
\end{center}
\begin{quote}
Außerdem wurde die Summe der Produkte aus den Quadraten der Fehler und den Gewichten,
für die im Original zweimal $9{,}6184$ angegeben ist, in $9{,}6154$ geändert.
\end{quote}
Die folgende Angabe blieb unverändert.
\par\noindent\begin{tabular}{@{}p{.065\linewidth}p{.91\linewidth}@{}}
S.\ 31&Bei Stern 9, Altona, das Gewicht des Kollimationsfehlers von $4{,}00$. Nach der Anzahl der
Beobachtungen auf S.\ 15 muss es $5{,}33$ betragen. Der Wert $4{,}00$ wurde beibehalten, weil mit ihm
das Endergebnis auf S.\ 32 abgeleitet wurde.
\end{tabular}

Ferner ist zu Art.\ 7, S.\ 35, Folgendes zu bemerken. Die Summe der Produkte aus den Quadraten
der in Art.\ 6 aufgeführten Unterschiede und der entsprechenden Anzahl der Beobachtungen ist dort
mit $292{,}8249$ angegeben. Da diese Summe den aus ihr abgeleiteten mittleren
Fehler einer Beobachtung von $1''{,}5308$ nicht genau ergab, wurden die Werte nachgerechnet.
Die Nachrechnung ergab anstelle des ersten Werts $291{,}8422$ und damit als mittleren Fehler
einer Beobachtung $\pm1''{,}5280$ sowie als mittleren Fehler des Breitenunterschieds
$\pm0''{,}1036$ anstelle der $\pm0''{,}1038$ des Originals.

Schließlich wurde im IV.\ Kapitel in der Tabelle auf S.\ 53 bei Stern 31 der im Original
angegebene Wert der Zenitdistanz von $+4^\circ48'58''{,}28$ in $+4^\circ57'58''{,}28$ geändert.
Der letztere Wert ergibt sich rückwärts aus der nebenstehenden Breite und der auf S.\ 9
aufgeführten Deklination des Sterns.

Die von \textsc{Gauss} in der Anzeige der Bestimmung des Breitenunterschieds etc., S.\ 61,
erwähnte Abhandlung zur lappländischen Gradmessung ist wahrscheinlich die von
O.\ A.\ \textsc{Rosenberger}, Über die, auf Veranstaltung der französischen Academie, während
der Jahre 1736 und 1737 in Schweden vorgenommene Gradmessung. (Astr.\ Nachr.\ Sechster Band,
1828. S.\ 1--32).
\par\hfill\textsc{Krüger, Börsch.}\par
}
\begin{center}\rule{.25\textwidth}{.4pt}\end{center}

\clearpage
% Unnumbered section title, principal PDF 72 (position 65).
\vspace*{45mm}
\begin{center}
{\Huge ERDELLIPSOID}\par\vspace{10mm}
{\large UND}\par\vspace{10mm}
{\Huge GEODÄTISCHE LINIE.}
\end{center}
\clearpage
% Blank verso, principal PDF 73 (position 66).
\null
\clearpage
% Unnumbered Nachlass opening, principal PDF 74 (position 67).
\begin{center}
{\Large\textls[180]{NACHLASS.}}\par\vspace{12mm}
[1.]\par
Das [Erd]ellipsoid.
\end{center}
Es sei $a$ die große, $b$ die kleine Halbachse. Für einen Ort sei $\varphi$ die astronomische
Polhöhe, $\varphi'$ die sogenannte verbesserte Polhöhe [und $\psi$ die reduzierte Polhöhe].
Der Abstand vom Mittelpunkt sei $r$, von der Erdachse $x$, vom Äquator $y$.

Da nach der Gleichung der Ellipse $\dfrac{xx}{aa}+\dfrac{yy}{bb}=1$ gilt, kann man
\[
x=r\cos\varphi'=a\cos\psi,\qquad y=r\sin\varphi'=b\sin\psi
\]
setzen. Bezeichnet man mit $s$ den elliptischen Bogen zwischen dem Ort und dem Äquator, so gilt
\[
-\sin\varphi\cdot\mathrm ds=\mathrm dx,\qquad\cos\varphi\cdot\mathrm ds=\mathrm dy,
\]
also
\[
\sin\varphi\cdot\mathrm ds=a\sin\psi\cdot\mathrm d\psi,\qquad\cos\varphi\cdot\mathrm ds=b\cos\psi\cdot\mathrm d\psi.
\]
Hieraus folgt
\[
\operatorname{tang}\psi=\frac ba\operatorname{tang}\varphi,\qquad\operatorname{tang}\varphi'=\frac{bb}{aa}\operatorname{tang}\varphi.
\]
Ferner
\[
\begin{aligned}
\cos\psi&=\frac{a\cos\varphi}{\sqrt{(aa\cos\varphi^2+bb\sin\varphi^2)}}=\sqrt{\frac{\cos\varphi\cos\varphi'}{\cos(\varphi-\varphi')}},\quad&\cos\varphi'&=\frac{aa\cos\varphi}{\sqrt{(a^4\cos\varphi^2+b^4\sin\varphi^2)}}\\
\sin\psi&=\frac{b\sin\varphi}{\sqrt{(aa\cos\varphi^2+bb\sin\varphi^2)}}=\sqrt{\frac{\sin\varphi\sin\varphi'}{\cos(\varphi-\varphi')}},&\sin\varphi'&=\frac{bb\sin\varphi}{\sqrt{(a^4\cos\varphi^2+b^4\sin\varphi^2)}}
\end{aligned}
\]
\[
r=\sqrt{\frac{a^4\cos\varphi^2+b^4\sin\varphi^2}{aa\cos\varphi^2+bb\sin\varphi^2}}=a\sqrt{\frac{\cos\varphi}{\cos\varphi'\cos(\varphi-\varphi')}}=b\sqrt{\frac{\sin\varphi}{\sin\varphi'\cos(\varphi-\varphi')}}
\]
\noindent\hfill $9^*$\par

\clearpage
\noindent 68\hfill NACHLASS.\hfill\mbox{}\par
\[
\frac{a\,\mathrm d\psi}{\cos\psi^2}=\frac{b\,\mathrm d\varphi}{\cos\varphi^2},\quad\text{also}\quad\frac{\mathrm ds}{\mathrm d\varphi}=\frac{bb\cos\psi^3}{a\cos\varphi^3}=\frac{aa\sin\psi^3}{b\sin\varphi^3}=\frac{aabb}{(aa\cos\varphi^2+bb\sin\varphi^2)^{\frac32}}
\]
\[
ee=1-\frac{bb}{aa}.
\]
Ist ein Quadrant des Erdmeridians $=Q$, so ist ein Quadrant des Äquators [vgl.\ Band IV, S.\ 330]
\[
\begin{aligned}
&=Q\cdot(1+\tfrac14ee+\tfrac7{64}e^4+\tfrac{15}{256}e^6+\tfrac{579}{16384}e^8+\tfrac{1515}{65536}e^{10}+\ldots)\\
&=Q:(1-\tfrac14ee-\tfrac3{64}e^4-\tfrac5{256}e^6-\tfrac{175}{16384}e^8-\tfrac{441}{65536}e^{10}-\ldots).
\end{aligned}
\]
\begin{center}\rule{.25\textwidth}{.4pt}\end{center}
Für die Abplattung $\tfrac1{310}$ ist $e=\dfrac{\sqrt{619}}{310}$, $\log e=8{,}9044836$
\[
\begin{array}{l@{\quad}l@{\quad}l}
ee\ldots\ 7{,}808\,9672&e^4\ldots\,\ldotp\ 5{,}617\,9344&e^6\ldots\,\ldotp\,\ldotp\ 3{,}426\,9016\\[4pt]
\tfrac14\ldots\,\ldotp\ 9{,}397\,9400&\tfrac7{64}\ldots\ 9{,}038\,9180&\tfrac{15}{256}\ldots\ 8{,}767\,8513
\end{array}
\]
Quadr.\ des Äqu.\ [wenn $Q=10\,000\,000$ Meter angesetzt wird:]
\[
\begin{array}{r@{\quad}c@{\quad}r}
10\,016\,103^{\mathrm m}{,}015&&\\
45{,}379&&\\
0{,}157&&\\\cline{1-1}
10\,016\,148^{\mathrm m}{,}551&\ldotp\,\ldotp\,\ldotp\,\ldotp\,\ldotp\,\ldotp\,\ldotp&7{,}000\,7008\\
&\left[\dfrac\pi2\right]\ldots&0{,}196\,1199\\\cline{3-3}
&\log a=&6{,}804\,5809.
\end{array}
\]
\begin{center}\rule{.25\textwidth}{.4pt}\end{center}
\[
\text{Abplattung}=\omega,\qquad b=a(1-\omega),\qquad ee=(2-\omega)\omega.
\]
\[
e=\sin E,\qquad\frac{1-\sqrt{(1-ee)}}{1+\sqrt{(1-ee)}}=\operatorname{tang}\tfrac12E^2.
\]
Für $\omega=\dfrac1{302{,}68}$ [ist]
\[
E=4^\circ39'31''{,}355,\qquad\operatorname{tang}\tfrac12E\ldots\ 8{,}609\,3534
\]
\[
\begin{array}{l@{\quad}l@{\quad}l}
ee\ldots\ 7{,}819\,3287&e^4\ldots\,\ldotp\ 5{,}638\,6574&e^6\ldots\,\ldotp\ 3{,}457\,9861\\[4pt]
\tfrac14\ldots\ 9{,}397\,9400&\tfrac7{64}\ldots\ 9{,}038\,9180&\tfrac{15}{256}\ldotp\,\ldotp\ 8{,}767\,8513
\end{array}
\]
[Für $Q=10000000$ Meter ist mithin der Quadrant des Äquators:]

\clearpage
\noindent\mbox{}\hfill DAS ERDELLIPSOID.\hfill 69\par
\[
\begin{array}{r@{\quad}c@{\quad}r}
10\,016\,491^{\mathrm m}{,}8250&&\\
47{,}5954&&\\
0{,}1682&&\\\cline{1-1}
10\,016\,539^{\mathrm m}{,}5886&\ldotp\,\ldotp\,\ldotp\,\ldotp\,\ldotp\,\ldotp\,\ldotp\,\ldotp\,\ldotp\,\ldotp\,\ldotp\,\ldotp\,\ldotp\,\ldotp&7{,}000\,7177\,118\\
&\left[\dfrac\pi2\right]\ldots&0{,}196\,1198\,770\\\cline{3-3}
\multicolumn{2}{r}{\log a\text{ (in Metern)}=}&6{,}804\,5978\,348
\end{array}
\]
[Da der] Meridianquadrant $=5130878{,}3$ Toisen [ist, so ist]
\[
\log a\text{ (in Toisen)}=6{,}514\,7893\,106.
\]
\begin{center}\rule{.25\textwidth}{.4pt}\end{center}
Zahlen, das Erdsphäroid betreffend.
\[
\text{Abplattung}=\frac1{302{,}78},\quad ee=\frac{604{,}56}{(302{,}78)^2},\quad\log ee=7{,}81918\,50399{\scriptstyle45}.
\]
[Der Umfang des Meridians ist $=$] Peripherie des Äquators mal
\[
\begin{aligned}
\bigg(1-\tfrac14ee&-\frac{1.3}{4.16}e^4-\frac{1.3.15}{4.16.36}e^6-\frac{1.3.15.35}{4.16.36.64}e^8-\frac{1.3.15.35.63}{4.16.36.64.100}e^{10}\\
&\bigg[-\frac{1.3.15.35.63.99}{4.16.36.64.100.144}e^{12}-\frac{1.3.15.35.63.99.143}{4.16.36.64.100.144.196}e^{14}\bigg]-\ldots\bigg).
\end{aligned}
\]
\[
\begin{array}{rcl}
\tfrac14ee&=&\multicolumn{1}{r}{0{,}00164\,86370\,23537\,4099}\\[3pt]
\tfrac3{64}e^4&=&\multicolumn{1}{r}{20385\,03026\,5312}\\[3pt]
\tfrac5{256}e^6&=&\multicolumn{1}{r}{56\,01252\,6094}\\[3pt]
\tfrac{175}{16384}e^8&=&\multicolumn{1}{r}{20200\,3209}\\[3pt]
\tfrac{441}{65536}e^{10}&=&\multicolumn{1}{r}{83\,9236}\\[3pt]
\tfrac{4851}{1048576}e^{12}&=&\multicolumn{1}{r}{3805}\\[3pt]
\left[\tfrac{693693}{205520896}e^{14}\right]&=&\multicolumn{1}{r}{18}\\\cline{3-3}
&&0{,}00165\,06811\,48101\,1773\\[3pt]
\text{Summe der Reihe}&=&0{,}99834\,93188\,51898\,8227
\end{array}
\]
\[
\begin{array}{rcl}
\text{Logarithm.}&\ldots&9{,}99928\,25259{\scriptstyle92}\\
\dfrac{20\,000\,000}{\pi}&\ldotp\,\ldotp\,\ldotp\,\ldotp\,\ldotp\,\ldotp&6{,}80388\,01229{\scriptstyle70}\\\cline{3-3}
\text{\raisebox{.5ex}{.}}\quad\log a&=&6{,}80459\,75969{\scriptstyle78}.
\end{array}
\]

\clearpage
\noindent 70\hfill NACHLASS.\hfill\mbox{}\par
\begin{center}[2.]\par
[Gleichung der Vertikalebene des Rotationsellipsoids.]\end{center}
Für die Orte im Schnitt der Oberfläche des Ellipsoids mit einer Ebene, die bei der Länge $0$
und der Polhöhe $\varphi^0$ senkrecht auf der Oberfläche steht und den nördlichen Teil des
dortigen Meridians unter dem Winkel $\zeta^0$ schneidet, lautet die Bedingungsgleichung
\[
\frac{\cos\varphi\cos\lambda\sin\varphi^0}{\sqrt{(1-ee\sin\varphi^2)}}+\frac{\cos\varphi\sin\lambda\operatorname{cotang}\zeta^0}{\sqrt{(1-ee\sin\varphi^2)}}-\frac{(1-ee)\sin\varphi\cos\varphi^0}{\sqrt{(1-ee\sin\varphi^2)}}=\frac{ee\cos\varphi^0\sin\varphi^0}{\sqrt{(1-ee\sin\varphi^{0^2})}}
\]
oder einfacher
\[
\cos\varphi\cos\lambda\sin\varphi^0+\cos\varphi\sin\lambda\operatorname{cotang}\zeta^0-(1-ee)\sin\varphi\cos\varphi^0=ee\cos\varphi^0\sin\varphi^0\sqrt{\frac{1-ee\sin\varphi^2}{1-ee\sin\varphi^{0^2}}}.
\]
\begin{center}[3.]\end{center}
Gleichung des [Rotations]ellipsoids in Bezug auf eine Tangentialebene:
\[
xx(1-ee\sin\varphi^2)+yy(1-ee)+xzee\sin2\varphi+zz(1-ee\cos\varphi^2)-\frac{2bbz}{a\sqrt{(1-ee\sin\varphi^2)}}=0.
\]
[Die $x$-Achse ist Tangente an die Meridianellipse und nach Süden positiv. Die $z$-Achse
fällt mit der Richtung der Normalen zusammen und ist zum Inneren des Ellipsoids hin positiv.]

[Man hat auch:]
\[
\left(\frac{x\sin\varphi-z\cos\varphi}{a}\right)^2+\left(\frac{x\cos\varphi+z\sin\varphi}{b}\right)^2+\frac{yy}{aa}=\frac{2z}{a\sqrt{(1-ee\sin\varphi^2)}}.
\]
[Setzt man]
\[
\begin{aligned}
x&=s\cos\delta\cos\theta\\
y&=s\cos\delta\sin\theta\\
z&=s\sin\delta,
\end{aligned}
\]
[wobei $\delta$ der Depressionswinkel und $\theta$ das Azimut ist, so ergibt sich]
\[
s\left(1+\frac{ee}{1-ee}(\cos\varphi\cos\theta\cos\delta+\sin\varphi\sin\delta)\right)^2\sqrt{(1-ee\sin\varphi^2)}=2a\sin\delta.
\]
[Näherungsweise gilt]
\[
\sin\lambda=\frac{s\cos\delta\sin\theta\sqrt{(1-ee\sin\varphi^2)}}{a\cos\varphi}.
\]
\begin{center}\rule{.18\textwidth}{.4pt}\end{center}

\clearpage
\noindent\mbox{}\hfill DAS ERDELLIPSOID.\hfill 71\par
\vspace{12mm}
\begin{center}\textbf{BEMERKUNGEN.}\end{center}
{\small
Die Notiz [1] stammt aus einem Notizbuch. Die zugehörigen Zahlenwerte fanden sich in zwei
anderen Notizbüchern. Die Zahlen zum Ellipsoid sind auf der letzten Seite von \textsc{Gauss}' Exemplar
der »Mathematischen Abhandlungen von Dr.\ H.\ F.\ \textsc{Scherk}. Berlin, 1825« eingetragen.
Der Abplattungswert $\tfrac1{310}$ war der Wert, den General von \textsc{Müffling} den preußischen
Vermessungen zugrunde legte. Wie aus dem Brief an \textsc{Olbers} vom 18.\ April 1822 hervorgeht,
hat \textsc{Gauss} zu Beginn seiner Gradmessungsarbeiten aus Versehen anstelle des Werts von
\textsc{Walbeck} die Abplattung $\dfrac1{302{,}68}$ verwendet. Tatsächlich finden sich mehrere
Tabellen und Rechnungen, die darauf beruhen. \textsc{Gauss} wollte diesen Wert zunächst
beibehalten, entschloss sich später aber (siehe den Brief von \textsc{Gauss} an \textsc{Olbers}
vom 1.\ März 1827), zu \textsc{Walbeck}s Wert $\dfrac1{302{,}78}$ überzugehen. Dieser liegt
also der Berechnung der Messungen von \textsc{Gauss} zugrunde. Die Länge des Meridianquadranten
beträgt nach \textsc{Walbeck} $5130878{,}22$ Toisen (Dissertatio de forma et magnitudine telluris
ex dimensis arcubus meridiani definiendis. Aboae, 1819).

[2] und der erste Teil von [3] stammen aus demselben Notizbuch wie [1]. Der zweite Teil von
[3] fand sich auf einem einzelnen Blatt.
\par\hfill\textsc{Krüger, Börsch.}\par
}
\begin{center}\rule{.25\textwidth}{.4pt}\end{center}
